\section{食品表示：どの情報が選択と安全を支えるか}
\label{theory:compliance:foodlabel}
期限、原産国・原料原産地、成分・アレルギー、事業者情報は、それぞれ別の判断を支える。
賞味期限は定めた保存条件のもとで期待される品質が保たれる期間、
消費期限は安全性を欠くおそれがないと認められる期間である。
いずれも保存条件が前提で、開封後や不適切な温度に置いた食品の無条件保証ではない。
賞味期限を過ぎたら直ちに全て食べられない、という安全期限との混同を避ける。\cite{saf53_food_expiry,saf53_food_loss2013}

輸入加工食品の最終加工国と、国内で加工した食品に使う原料の産地を分ける。
平成25年度の品目・原料の条件を、後年の広い対象へ変えない。
食品表示法の公布2013年と施行2015年も区別する。
アレルギー表示の指定は見直されるため、当時の指定品目を現在の数に書き換えない。\cite{saf53_food_notice}
表示は選択だけでなく、事故時の対象製品と責任を追跡する情報でもある。

\section{不実証広告：表示を裏付ける二つの条件}
\label{theory:compliance:adevidence}
効果を示す資料があるという事実だけでは、広告の主張は確かめられない。
資料自体の客観性と、表示された対象・条件・効果への対応の両方が必要である。
別の材料の試験や特殊条件の結果を、製品の一般的な利用へそのまま拡張しない。
行政が資料を求めても自由に断れるという制度ではなく、所定の不提出には制度上の効果がある。\cite{saf53_ad_overview}

自社試験でも妥当な方法で実証できるが、実施者が専門家であるだけでは足りない。
専門分野で一般に認められない特異な見解や、新しい技術を称することによって
試験・調査の実証を免れるものではない。
経験則を根拠とする場合も、存在を裏付ける専門的な見解等を確認する。\cite{saf53_ad2016}
2016年の説明では資料提出は原則15日で、新しい試験が必要という理由だけで延長できるとはしない。
2014年の設問へ現行条番号や後年の処分制度を補わず、当時の枠組みを保存する。
